\section{Discussion}
\label{sec:discussion}

\paragraph{What a soft-binding service should deploy.} No single soft binding serves every purpose, and the results argue for a layered design with a different lead family per modality. For images, a watermark and a fingerprint recover about the same share of attacked copies but different ones, so the union of their successes covers \val{fp:wm/image/DINOv2-S/either/pct}\% of attacked copies with TrustMark and DINOv2-S, against \val{fp:wm/image/DINOv2-S/wm/pct}\% and \val{fp:wm/image/DINOv2-S/fp/pct}\% for either alone (Section~\ref{sec:together}). Among openly licensed watermarks, PixelSeal at a strength just above its robustness threshold is the best balance of perceptibility, robustness and speed; its margin at 0.15 is small on small images, so a deployment expecting heavier compression should move towards the default of 0.2 (Section~\ref{sec:wm-res-image}). Among openly licensed fingerprints, a DINOv2 embedding recovers the most copies and is the only one that survives large crops and rotations; its keyed 256-bit projection costs a few points of detection, and stores in the same 32 bytes as PDQ. In the evaluated inversion experiment, a decoder trained under a different key did not identify the source above chance (Section~\ref{sec:fp-res-security}); this does not establish that no attacker can invert it. The registered open standard, the ISCC Image-Code, is precise but loses roughly half of the transformed copies. For audio, fingerprints cover more queries: Chromaprint and audfprint recover far more attacked copies than the 16-bit AudioSeal key, whose base model is nevertheless the best watermark operating point. For video, expected-key watermark verification covers more queries: in the combination experiment Video Seal passed that test after the crops, flips and insertions that break every open video fingerprint, and in the watermarking track PixelSeal at 0.4 was the best-balanced video watermark. The strongest image copy detectors, SSCD and the ISC21 winner, would add a further 6--10 points but cannot be shipped as released, because their weights were trained on data licensed for non-commercial research.

\paragraph{The decision rule decides the false positives.} For both families the false-match rate belongs to the complete decision rule, not to the model. TrustMark's error-correcting decoder reports a valid codeword on several percent of unmarked images, yet none of the same decodes reached the expected-payload threshold on any of the \val{wm:fpr/image/TrustMark-Q/own_src_n} unmarked images (Section~\ref{sec:wm-res-fpr}); a provenance verifier resolves a soft binding against a manifest and can apply that test. The combination experiment shows the same for keys decoded from attacked, never-registered content (Section~\ref{sec:together}): a decoded key must be checked against the registry exactly as a fingerprint score is thresholded. Payload size limits what can be shown. With 16 bits, even an exact match leaves a chance of $1.5\times10^{-5}$ per key, so the 16-bit audio models rest on their blind detectors, whose observed bound over \val{wm:fpr/audio/WavMark/blind_src_n} clips does not establish $10^{-3}$. For fingerprints, $Np$ is the expected number of false matches per negative query for $N$ references and average pair-level false-match probability $p$, and bounds the probability of at least one false match from above. A pair-level rate at most $10^{-9}$ would therefore suffice for a 1\% query-level bound at $10^{7}$ assets; this operating point is below the empirically resolved range here. We expected a second, geometric verification stage to supply the missing precision, and measured it; under a matched calibration false-binding constraint it did not raise detection in the evaluated image pipelines, because the remaining false matches are near-copies of registered designs that also match geometrically (Section~\ref{sec:fp-res-fpr}). What alignment against the stored original does provide is localization of an edit, most of which survives when the original must first be retrieved.

\paragraph{Catalogues are full of near-copies.} The largest effect in the fingerprinting track is not the choice of method but the treatment of variants: when recoloured or re-mounted versions of one design count as distinct assets, every method's calibrated detection falls to 0.24 or less. Learned descriptors bind such variants far more readily than hashes, which is correct if the question is ``is this the same artwork'' and wrong if it is ``is this the asset this manifest was signed for''. A soft-binding service must decide which question it answers, group variants of a design under one registry entry if it answers the first, and treat a match as a pointer to candidate manifests rather than as an identification.

\paragraph{Neither family is a security boundary.} White-box attacks evaded every learned fingerprint with perturbations of a few grey levels, surrogate gradients evaded the DCT hashes, targeted forgeries bound unrelated images to registered assets, and pitch, tempo or small crops broke the audio and video fingerprints. Regeneration and adversarial attacks are known to remove or evade pixel-level watermarks \cite{zhao2024removable, jiang2023evading}; our watermarking track measured routine processing only. A soft-binding match should therefore be presented as a pointer to a candidate manifest that must be confirmed, never as evidence that content is authentic. The value written into a \texttt{c2pa.soft-binding} assertion is disclosed to everyone who reads the manifest, and our inversion decoder recovered recognisable thumbnails from unkeyed descriptors and hashes, so services should keep raw descriptors server-side and publish only keyed or opaque identifiers.

\paragraph{Fingerprint the published asset.} Embedding a watermark is itself a transformation. It moved the ISCC Image-Code of \val{fp:wm/image/ISCC-Image-64/drift/pct}\% of images past its own match threshold, and a registry holding the fingerprint of the unmarked original lost matches accordingly (Section~\ref{sec:together}). A service that does both must compute the fingerprint on the watermarked asset it publishes.

\paragraph{Reproducibility across hardware and runs.} Provenance claims should not depend on which machine checks them. A watermark decoded on an x86 host lost payload bits on an arm64 host; given identical RGB arrays both hosts agreed to within floating-point noise, so the cause is the platform-specific conversion from YUV to RGB that feeds the detector (Section~\ref{sec:wm-res-stability}). Converting colour in a fixed, portable implementation, or running detectors on YUV input, would remove the effect. Two classical image hashes, Marr--Hildreth and wHash, were likewise not bit-identical across architectures. Reproducing a benchmark also needs more pinned than its corpus: our two runs of the audio fingerprint track drew different background music for the babble transformations, because the pool came from incidental surplus downloads, and the results moved accordingly until the pool was pinned.

\paragraph{Limitations.} The watermark corpus is small by the standards of model training (85 images, 138 audio inputs from \val{wm:audio/n_sources} clips, six 2.5-s video clips); it separates the methods on the gate and on false positives, but the video comparisons are descriptive, FLIP and ColorVideoVDP were computed on subsets, and PESQ is reported on speech only. ABO, the fingerprint image corpus, is a product-photo catalogue whose white backgrounds and repeated designs make some things harder and others easier than natural photographs, as the comparison with DISC21 shows. The fingerprint audio negatives are music only. The video calibration set contains \val{fp:count/video/neg_cal} queries from \val{fp:video/cal_sources} negative clips and supports reporting the query-level 1\% and pair-level $10^{-4}$ targets under our resolution rule, with limited precision. The achieved held-out query-level false-match rates exceed the 1\% calibration target. Source-clip resampling accounts for transformations of the same clip, but does not by itself account for dependence between clips from the same film. The security experiments used small samples and simple attacks, so they bound attack difficulty from above. Latency was measured on one GPU type, and cost estimates use on-demand prices at one date and assume full utilisation. Strength settings were chosen around each method's default and do not trace a full rate--distortion curve. The licence status of several models changed during 2024--2026; we report the status recorded at the time of measurement.
